\begin{proposition}
\label{proposition-left-derived-exists}
Let $F : \mathcal{A} \to \mathcal{B}$ be a right exact functor
of abelian categories. Let $\mathcal{P} \subset \Ob(\mathcal{A})$ be a
subset. Assume
\begin{enumerate}
\item $\mathcal{P}$ contains $0$, is closed under (finite) direct sums,
and every object of $\mathcal{A}$ is a quotient of an
element of $\mathcal{P}$,
\item for any bounded above acyclic complex $P^\bullet$ of
$\mathcal{A}$ with $P^n \in \mathcal{P}$ for all $n$ the
complex $F(P^\bullet)$ is exact,
\item $\mathcal{A}$ and $\mathcal{B}$ have colimits
of systems over $\mathbf{N}$,
\item colimits over $\mathbf{N}$ are exact in both
$\mathcal{A}$ and $\mathcal{B}$, and
\item $F$ commutes with colimits over $\mathbf{N}$.
\end{enumerate}
Then $LF$ is defined on all of $D(\mathcal{A})$.
\end{proposition}

\begin{proof}
By (1) and Lemma \ref{lemma-subcategory-left-resolution} for any bounded
above complex $K^\bullet$ there exists a quasi-isomorphism
$P^\bullet \to K^\bullet$ with $P^\bullet$ bounded above and
$P^n \in \mathcal{P}$ for all $n$. Suppose that
$s : P^\bullet \to (P')^\bullet$ is a quasi-isomorphism of bounded
above complexes consisting of objects of $\mathcal{P}$. Then
$F(P^\bullet) \to F((P')^\bullet)$ is a quasi-isomorphism because
$F(C(s)^\bullet)$ is acyclic by assumption (2). This already shows that
$LF$ is defined on $D^{-}(\mathcal{A})$ and that a bounded above
complex consisting of objects of $\mathcal{P}$ computes $LF$, see
Lemma \ref{lemma-find-existence-computes}.

\medskip\noindent
Next, let $K^\bullet$ be an arbitrary complex of $\mathcal{A}$.
Choose a diagram
$$
\xymatrix{
P_1^\bullet \ar[d] \ar[r] & P_2^\bullet \ar[d] \ar[r] & \ldots \\
\tau_{\leq 1}K^\bullet \ar[r] & \tau_{\leq 2}K^\bullet \ar[r] & \ldots
}
$$
as in Lemma \ref{lemma-special-direct-system}. Note that
the map $\colim P_n^\bullet \to K^\bullet$ is a quasi-isomorphism
because colimits over $\mathbf{N}$ in $\mathcal{A}$ are exact
and $H^i(P_n^\bullet) = H^i(K^\bullet)$ for $n > i$. We claim that
$$
F(\colim P_n^\bullet) = \colim F(P_n^\bullet)
$$
(termwise colimits) is $LF(K^\bullet)$, i.e., that $\colim P_n^\bullet$
computes $LF$. To see this, by Lemma \ref{lemma-find-existence-computes},
it suffices to prove the following claim. Suppose that
$$
\colim Q_n^\bullet = Q^\bullet
\xrightarrow{\ \alpha\ }
P^\bullet = \colim P_n^\bullet
$$
is a quasi-isomorphism of complexes, such that each
$P_n^\bullet$, $Q_n^\bullet$ is a bounded above complex whose terms are
in $\mathcal{P}$ and the maps $P_n^\bullet \to \tau_{\leq n}P^\bullet$ and
$Q_n^\bullet \to \tau_{\leq n}Q^\bullet$ are quasi-isomorphisms.
Claim: $F(\alpha)$ is a quasi-isomorphism.

\medskip\noindent
The problem is that we do not assume that $\alpha$ is given as a colimit
of maps between the complexes $P_n^\bullet$ and $Q_n^\bullet$. However,
for each $n$ we know that the solid arrows in the diagram
$$
\xymatrix{
& R^\bullet \ar@{..>}[d] \\
P_n^\bullet \ar[d] &
L^\bullet \ar@{..>}[l] \ar@{..>}[r] &
Q_n^\bullet \ar[d] \\
\tau_{\leq n}P^\bullet \ar[rr]^{\tau_{\leq n}\alpha} & &
\tau_{\leq n}Q^\bullet
}
$$
are quasi-isomorphisms. Because quasi-isomorphisms form a multiplicative
system in $K(\mathcal{A})$ (see Lemma \ref{lemma-acyclic})
we can find a quasi-isomorphism
$L^\bullet \to P_n^\bullet$ and map of complexes $L^\bullet \to Q_n^\bullet$
such that the diagram above commutes up to homotopy. Then
$\tau_{\leq n}L^\bullet \to L^\bullet$ is a quasi-isomorphism.
Hence (by the first part of the proof) we can find a bounded above
complex $R^\bullet$ whose terms are in $\mathcal{P}$ and a quasi-isomorphism
$R^\bullet \to L^\bullet$ (as indicated in the diagram). Using the result
of the first paragraph of the proof we see that
$F(R^\bullet) \to F(P_n^\bullet)$ and $F(R^\bullet) \to F(Q_n^\bullet)$
are quasi-isomorphisms. Thus we obtain an isomorphism
$H^i(F(P_n^\bullet)) \to H^i(F(Q_n^\bullet))$ fitting into the commutative
diagram
$$
\xymatrix{
H^i(F(P_n^\bullet)) \ar[r] \ar[d] &
H^i(F(Q_n^\bullet)) \ar[d] \\
H^i(F(P^\bullet)) \ar[r] &
H^i(F(Q^\bullet))
}
$$
The exact same argument shows that these maps are also compatible
as $n$ varies. Since by (4) and (5) we have
$$
H^i(F(P^\bullet)) =
H^i(F(\colim P_n^\bullet)) =
H^i(\colim F(P_n^\bullet)) = \colim H^i(F(P_n^\bullet))
$$
and similarly for $Q^\bullet$ we conclude that
$H^i(\alpha) : H^i(F(P^\bullet) \to H^i(F(Q^\bullet)$ is an isomorphism
and the claim follows.
\end{proof}